\unnumberedchapter{Conclusion générale}

Dans le cadre de ce projet, une démarche d'audit de sécurité et de durcissement a été menée sur deux actifs du système d'information : un poste de travail Windows 11 et un pare-feu FortiGate. L'objectif était d'évaluer leur niveau de conformité à partir des référentiels de sécurité retenus, d'identifier les écarts et de mettre en œuvre des mesures correctives adaptées.

Les audits initiaux ont permis d'identifier plusieurs paramètres nécessitant des améliorations. Une phase de durcissement a ensuite été réalisée en combinant des mesures automatisées à l'aide de scripts et des configurations effectuées manuellement. Des simulations de sécurité ont également été réalisées avant et après durcissement afin de vérifier l'efficacité des mesures appliquées.

Les résultats obtenus montrent une amélioration significative du niveau de conformité. Pour Windows 11, le taux de conformité est passé de \textbf{25\,\% à 85\,\%}, tandis que celui du FortiGate est passé de \textbf{8,9\,\% à 80,4\,\%}. Ces résultats montrent l'intérêt d'une approche structurée combinant audit, remédiation et validation.

\section*{Perspectives}

La principale perspective de ce travail consiste à appliquer la démarche proposée dans un \textbf{environnement réel de production}. Une telle mise en œuvre nécessiterait d'adapter les mesures de durcissement à l'architecture existante, aux contraintes opérationnelles et aux politiques de sécurité de l'organisation.

À terme, cette démarche pourrait être étendue à un plus grand nombre d'actifs du système d'information et intégrée dans un processus régulier d'audit et de contrôle de conformité. La centralisation des journaux, l'automatisation des contrôles, la surveillance continue ainsi que la réalisation périodique de tests de sécurité constitueraient également des axes d'amélioration.

Ainsi, ce projet constitue une base pour une démarche de \textbf{sécurisation continue des actifs du système d'information}, depuis l'identification des écarts jusqu'à leur correction et leur validation.
